\section{Conclusion and Future Work}

This thesis set out to support the migration of existing Elixir codebases to the language's new set-theoretic type system along the two fronts on which type information is available: translating the types that are already declared as Erlang-style TypeSpecs, and inferring those that are not.
Chapter~\ref{ch:formalisation} formalised the translation as a set of equations from TypeSpec syntax to Elixir Types, isolating the constructs that cannot be translated precisely and approximating them with the gradual type \texttt{dynamic()}.
Chapter~\ref{ch:implementation} realised that translation as a prototype, applied it across 383 repositories (drawn from 54 distinct owners) to produce \num{23073} translated annotations, and cross-validated each against both Dialyzer and the new typechecker.

The agreement study established that the new typechecker is the substantially stricter of the two analyses, rejecting \num{6011} annotations against Dialyzer's 285, while subsuming the majority of Dialyzer's own findings (a subsumption rate of 60.7\%).
Where both tools resolve the same types the typechecker reproduces Dialyzer's warnings; the residual cases are attributable not to disagreement but to whole-program analyses the typechecker does not yet perform, to \texttt{@opaque}-enforcement it does not model, and to translation-time \texttt{dynamic()} approximation, so the two analyses are complementary rather than nested.
The cross-validation yielded \num{7498} annotations accepted by both tools; these constitute the checker-verified supervision signal for the prediction task.

The prediction study fine-tuned a code LLM to recover these annotations and, through a two-track design, asked whether gradual types should be admitted into the supervision signal.
The result is a trade-off rather than a single best choice.
A model trained on the precise, dynamic-free pool and one trained on the larger gradual-inclusive pool invert across two equally legitimate notions of success: the gradual-inclusive model recovers more typechecker-accepted annotations (raw pass rate 75.0\% versus 63.7\%), which serves a \emph{low-noise} migration whose immediate aim is to make a codebase compile; the precise model produces more annotations that are accepted \emph{and} genuinely static (safe-and-precise rate 63.7\% versus 29.7\%) and never resorts to the gradual type, which serves type safety and downstream usefulness.
Critically, the gradual-inclusive model confirms the concern that motivated this work, namely that a model can learn to over-emit \texttt{dynamic()} to mute type errors, lifting the bare pass rate while halving precision; it does so despite training on more than twice the data.
The practical guidance is therefore a tunable trade-off rather than a single recommended model: admit gradual types when a low-noise transition is the priority, and exclude them when precision is.

Taken together, the translation prototype and the two prediction models cover the two sources of type information in a codebase, and the checker-verified, two-track methodology offers a reusable template for building and evaluating type-annotation datasets for a young set-theoretic type system.
What remains is to compose the two methods into a single end-to-end migration of a complete project and to measure it directly, which we set out as the principal item of future work below.

\subsection{Limitations and Future Work}
\label{sec:future-work}

\paragraph{End-to-end total migration.}
This thesis defines \emph{total migration} as the state in which every function in a project carries a typechecker-accepted Elixir Type, and it establishes the two enablers separately: a translation prototype that converts existing \texttt{@spec}s (Chapter~\ref{ch:implementation}), and an LLM that predicts annotations for the unspec'd majority (the two-track prediction study).
Composing the two into an end-to-end migration of a complete, previously-unannotated project is the natural next step and is left to future work.
Concretely, such a study would run the full pipeline, namely target discovery of untyped functions, annotation \emph{inference} (translation where a spec exists and prediction where it does not), \emph{injection} of the resulting annotation into the source, and recompilation in the original project context, and would report a \emph{migration rate}: the fraction of previously-untyped functions that become compile-clean and free of type warnings.
A case study over a handful of projects, with a per-function breakdown of outcomes (clean compile, type warning, or compile failure from a syntactic or unresolved-reference error), would connect the compatibility metric of Subsection~\ref{sec:compatibility-eval} with extrinsic migration success, testing directly whether compatibility with a reference predicts whether a prediction actually migrates.

\paragraph{Evaluation scope and dataset.}
Several limitations bound the present results and motivate the study above.
The extrinsic evaluation is conducted on \emph{held-out} projects rather than within a single evolving codebase.
The corpus itself is bounded by what the custom compiler's OTP-28 toolchain can build: projects whose dependencies do not compile under it never enter the dataset, so the results generalise to typecheckable Elixir code rather than to the language's full ecosystem.
The supervised dataset, while checker-verified, is modest in size (\num{7498} both-pass entries) and resolves no cross-project dependency types at training time, so the model never observes the remote types the translator was forced to widen to \texttt{dynamic()}, precisely the cases where prediction would most help.
Where the model does fail, the failures concentrate in framework code with non-obvious internal representations (packed-binary math types, opaque struct APIs) and in the recurring atom/integer-versus-binary confusion, suggesting that dependency-aware context at training time is the most promising lever for the next iteration.

\paragraph{Manual project preparation.}
Constructing and evaluating the corpus was not a fully automated pipeline: a number of projects required manual intervention before they would compile under the custom compiler.
The development-branch Elixir is built against a specific Erlang/OTP release (OTP-28), so projects pinning an older toolchain, or depending on libraries incompatible with it, did not build out of the box; bringing them up required pinning compatible dependency versions and adjusting build settings.
For instance, a project that aliased compilation to \texttt{--warnings-as-errors} turned the typechecker's own warnings into hard failures and had to be relaxed.
These interventions were applied uniformly and identically to both tracks, so they do not bias the comparison; but they mean the corpus cannot yet be regenerated end-to-end without human effort, and that scaling to a substantially larger project set would demand either more robust automated environment resolution or an accepted loss of coverage.

\paragraph{A compiler under active development.}
The set-theoretic type system targeted in this work ships in a development-branch compiler (Elixir~1.21.0-dev) whose type-checking behaviour is not yet stable, and the translation pipeline depends directly on it.
This coupling produced concrete incompatibilities between the translated annotations and the compiler.
The clearest is an arity mismatch in the \texttt{@assert\_type\_form} assertions the translation emits: for a function declaring specifications at several arities alongside a default-argument clause, the assertion generated for a higher arity could be associated by the compiler with a lower-arity clause, aborting compilation of the whole file; resolving these required repositioning the offending \texttt{@spec}s by hand.
Such mismatches are a property of a young, moving target rather than of the translation itself, so the empirical results reported here are tied to a particular snapshot of the compiler and can be expected to shift, most plausibly to improve, as its type checker and its handling of assertions mature.
